\clearpage
\subsection{Unendlich viele Lösungen parametrisieren}
\label{subsec:linear_systems_parametric_solution_sets}

Besitzt ein lösbares LGS unendlich viele Lösungen, wird seine Lösungsmenge
mithilfe von Parametern beschrieben.
Das folgende Beispiel zeigt dies anhand eines bereits bekannte LGS.

\begin{ExampleBox}{Eine Lösungsgerade parametrisieren}{linear_systems_parametric_solution}
  Dieses LGS wurde zuvor geometrisch betrachtet. Nun untersuchen wir es
  zusätzlich mithilfe des Rangs und lösen es mit dem Gauß-Verfahren:
  \[
    \begin{aligned}
      \textcolor{THMBlue}{x_1+x_2}&\textcolor{THMBlue}{=2},\\
      \textcolor{olive!80!black}{2x_1+2x_2}&\textcolor{olive!80!black}{=4}.
    \end{aligned}
    \quad\Longleftrightarrow\quad
    \left[
      \begin{array}{rr|r}
        \textcolor{THMBlue}{1}&\textcolor{THMBlue}{1}&\textcolor{THMBlue}{2}\\
        \textcolor{olive!80!black}{2}&\textcolor{olive!80!black}{2}&\textcolor{olive!80!black}{4}
      \end{array}
    \right].
  \]
  Die Vorwärtselimination liefert die Zeilenstufenform
  \[
    \left[
      \begin{array}{rr|r}
        \textcolor{THMBlue}{1}&\textcolor{THMBlue}{1}&\textcolor{THMBlue}{2}\\
        \textcolor{olive!80!black}{2}&\textcolor{olive!80!black}{2}&\textcolor{olive!80!black}{4}
      \end{array}
    \right]
    \quad
    \textcolor{black!55}{
      \left.\vphantom{\begin{array}{c}0\\0\end{array}}\right|
      \quad
      \begin{array}{l}
        \phantom{Z_2-2Z_1}\\
        Z_2-2Z_1
      \end{array}
    }
    \quad\Longrightarrow\quad
    \left[
      \begin{array}{rr|r}
        \textcolor{THMBlue}{1}&\textcolor{THMBlue}{1}&\textcolor{THMBlue}{2}\\
        \textcolor{olive!80!black}{0}&\textcolor{olive!80!black}{0}&\textcolor{olive!80!black}{0}
      \end{array}
    \right]
  \]
  Es gilt
  \[
    \left.
    \begin{aligned}
      n&=2,\\[0.5em]
      \operatorname{rang}(A)&=1,\\[0.5em]
      \operatorname{rang}([A\mid b])&=1
    \end{aligned}
    \quad\right\}
    \quad\Longrightarrow\quad
    \infty\text{-viele Lösungen}.
  \]

  \noindent
  \begin{minipage}[c]{0.62\linewidth}
    Um die Lösungsmenge geschickt anzugeben, führen wir einen Parameter
    $t\in\R$ ein und setzen
    \[
      x_2=\Cloze{1cm}{t}
      \quad\Longrightarrow\quad
      x_1=\Cloze{2cm}{2-t}.
    \]
    Damit lautet die Lösungsmenge
    \[
      \mathcal{L}
      =\Cloze{7.8cm}{
        \left\{
          \begin{pmatrix}2\\0\end{pmatrix}
          +t\begin{pmatrix}-1\\1\end{pmatrix}
          \;\middle|\;t\in\R
        \right\}
      }.
    \]
  \end{minipage}\hfill
  \begin{minipage}[c]{0.34\linewidth}
    \centering
    \CachedTikzFigure{content/resources/chapters/linear_systems/tikz/infinitely_many_solutions_lines.tex}
  \end{minipage}

  \smallskip
  Der Parameter $t$ wirkt wie ein Schieberegler: Für jeden Wert von $t$
  entsteht eine andere Lösung $(2-t,t)$. Durchläuft $t$ alle reellen Zahlen,
  durchläuft der zugehörige Punkt die gesamte Lösungsgerade $x_1+x_2=2$.
\end{ExampleBox}

\begin{KeyPointBox}{Anzahl der freien Variablen}{linear_systems_number_free_variables}
  Für ein lösbares LGS mit $n$ Unbekannten bestimmt die Differenz zwischen $n$
  und $\operatorname{rang}(A)$, wie viele unabhängige Parameter zur Beschreibung
  der Lösungsmenge benötigt werden. Man spricht hierbei von
  \emph{freien Variablen} \EnglishTerm{free variables}. Es gilt
  \[
    \boxed{
      \#\{\text{benötigte Parameter}\}
      =\#\{\text{freie Variablen}\}
      =n-\operatorname{rang}(A)
    }.
  \]
  \par\smallskip
  \hfill
  {\footnotesize\color{black!55}
    $\#\{\ldots\}\mathrel{\widehat{=}}\text{Anzahl der Elemente von }\{\ldots\}$.}
\end{KeyPointBox}
